\begin{document}

\section{Preliminaries}

\begin{definition}[Mathlib Preliminaries]
    \label{def:preliminaries}
    \lean{ContinuousOn}

    We use the following standard Mathlib notions throughout.
    \begin{itemize}
        \item \textbf{Continuity on a set} (\texttt{ContinuousOn}): A function $f : \alpha \to \beta$ between topological spaces is \emph{continuous on} $s \subseteq \alpha$ if for every $x \in s$, the filter map of $f$ along the neighborhood filter of $x$ restricted to $s$ tends to $f(x)$.
        \item \textbf{Derivative at a point} (\texttt{HasDerivAt}): For $F : \mathbb{F} \to \mathbb{F}$ over a nontrivially normed field, $F$ \emph{has derivative} $f'$ at $x$ if $F(x+h) = F(x) + f' \cdot h + o(h)$ as $h \to 0$.
        \item \textbf{Interval integral} (\texttt{intervalIntegral}): For $f : \mathbb{R} \to E$ and $a,b \in \mathbb{R}$, the integral $\int_a^b f(x)\,dx$ is defined as $\int_{[a,b]} f\,d\mu - \int_{[b,a]} f\,d\mu$, which is antisymmetric: $\int_a^b f = -\int_b^a f$.
    \end{itemize}
\end{definition}

\begin{definition}[Antiderivative]
    \label{def:antiderivative}
    \uses{def:preliminaries}
    \lean{FLT.FundamentalTheoremOfCalculus.IsAntiderivativeOn}

    A function $F : \mathbb{R} \to \mathbb{R}$ is an \emph{antiderivative} of $f$ on $(a,b)$ if for every $x \in (a,b)$, $F$ has derivative $f(x)$ at $x$ (in the sense of \texttt{HasDerivAt}). In the Lean formalization this is \texttt{IsAntiderivativeOn F f a b}, defined as $\forall x \in (a,b),\; \texttt{HasDerivAt}\; F\; (f\; x)\; x$.
\end{definition}

\section{Supporting Lemma}

\begin{lemma}[Zero derivative implies constant]
    \label{lem:zero_deriv_constant}
    \uses{def:preliminaries}
    \lean{FLT.FundamentalTheoremOfCalculus.zero_deriv_constant}

    Let $G : \mathbb{R} \to \mathbb{R}$ be continuous on $[a,b]$ and satisfy $\texttt{HasDerivAt}\; G\; 0\; x$ for all $x \in (a,b)$. Then $G(x) = G(a)$ for all $x \in [a,b]$.
\end{lemma}
\begin{proof}
    \uses{def:preliminaries}
    For any $x \in [a,b]$ with $a < x$, restrict $G$ to $[a,x]$. By the Mean Value Theorem (\texttt{exists\_hasDerivAt\_eq\_slope}), there exists $c \in (a,x)$ with $G'(c) = (G(x) - G(a))/(x - a)$. Since $G'(c) = 0$ by hypothesis, we get $G(x) = G(a)$.
\end{proof}

\section{Main Theorems}

\begin{theorem}[First Fundamental Theorem of Calculus]
    \label{thm:ftc1}
    \uses{def:preliminaries}
    \lean{FLT.FundamentalTheoremOfCalculus.ftc1}

    Let $f : \mathbb{R} \to \mathbb{R}$ be continuous on $[\min(a,b), \max(a,b)]$, and let $x \in (\min(a,b), \max(a,b))$. Then the function $F(u) = \int_a^u f(t)\,dt$ has derivative $f(x)$ at $x$. In Mathlib the core result is \texttt{intervalIntegral.integral\_hasDerivAt\_right}.
\end{theorem}
\begin{proof}
    \uses{def:preliminaries}
    Since $f$ is continuous on $[\min(a,b), \max(a,b)]$ and $x$ is in the interior, we establish three prerequisites for \texttt{intervalIntegral.integral\_hasDerivAt\_right}: (1) $f$ is interval integrable on $[a,x]$ via \texttt{ContinuousOn.intervalIntegrable}; (2) $f$ is strongly measurable at the neighborhood filter of $x$, using openness of the interior; (3) $f$ is continuous at $x$, since $x$ lies in the interior of the domain of continuity. The result then follows directly.
\end{proof}

\begin{theorem}[Second Fundamental Theorem of Calculus]
    \label{thm:ftc2}
    \uses{def:preliminaries}
    \lean{FLT.FundamentalTheoremOfCalculus.ftc2}

    Let $F : \mathbb{R} \to \mathbb{R}$ be continuous on $[\min(a,b), \max(a,b)]$ and have right derivative $f'(x)$ at every $x \in (\min(a,b), \max(a,b))$, i.e., $\texttt{HasDerivWithinAt}\; F\; (f'\; x)\; (x, +\infty)\; x$. If $f'$ is interval integrable on $[a,b]$, then
    \[
      \int_a^b f'(x)\,dx = F(b) - F(a).
    \]
    In Mathlib this is \texttt{intervalIntegral.integral\_eq\_sub\_of\_hasDeriv\_right}.
\end{theorem}
\begin{proof}
    \uses{def:preliminaries}
    Direct application of \texttt{intervalIntegral.integral\_eq\_sub\_of\_hasDeriv\_right}, given the continuity of $F$ on $[\min(a,b), \max(a,b)]$, the right-derivative hypothesis on the interior, and the integrability of $f'$ (which follows from \texttt{ContinuousOn.intervalIntegrable}).
\end{proof}

\end{document}
